\section{Introduction}
\label{sec:intro}

The volume of news produced online every day is staggering. Readers, researchers, and media organizations alike need ways to quickly understand what a given article is about and whether its tone is positive or negative. Doing this manually at scale is simply not feasible, which makes automated news analysis a practically important problem and a rich domain for applying NLP techniques.

Genre classification — assigning a topic category such as \textit{sport} or \textit{technology} to an article — is the most fundamental of these tasks. It underpins content recommendation systems, archive organization, and media monitoring. Beyond genre, two additional layers of information are particularly useful: (1) the specific terms and entities that distinguish an article from the background corpus, and (2) the overall sentiment of the coverage, which can reveal editorial tone or public mood around a topic.

This project builds an end-to-end pipeline, which we call \textbf{NewsScope}, that addresses all three tasks in sequence:

\begin{enumerate}
    \item \textbf{Genre classification:} A Multinomial Naive Bayes classifier assigns each article to one of five BBC categories.
    \item \textbf{Keyword and entity extraction:} TF-IDF scoring, POS-tag filtering, and named-entity recognition surface the most important terms and real-world entities in each article.
    \item \textbf{Sentiment analysis:} A Naive Bayes sentiment model trained on in-domain BBC data (pseudo-labeled by VADER) is compared against VADER's rule-based scores directly.
\end{enumerate}

The entire pipeline is implemented using NLTK \cite{bird2009natural} and scikit-learn, and demonstrated on live BBC RSS headlines.

% -----------------------------------------------------------------------
% TODO: Insert a figure showing the full pipeline architecture here.
%   Suggested figure: a block diagram with three parallel branches
%   (genre / keywords / sentiment) fed by a shared preprocessing step,
%   with a "Live RSS Demo" box at the bottom.
% -----------------------------------------------------------------------

The rest of this report is organized as follows. Section~\ref{sec:data} describes the dataset. Section~\ref{sec:method} details the preprocessing, feature extraction, and modeling choices. Section~\ref{sec:experiments} reports quantitative results and example outputs. Section~\ref{sec:discussion} discusses the findings and limitations. Section~\ref{sec:conclusion} concludes.

\section{Related Work}
\label{sec:related}

Text classification has a long history in NLP. Naive Bayes applied to bag-of-words representations was popularized by \citet{mccallum1998comparison}, who showed that even with its independence assumption, it achieves competitive results on document categorization benchmarks. Despite the rise of deep learning, it remains a strong baseline for topic classification due to its efficiency and interpretability — properties that make it well-suited for an educational NLP project.

TF-IDF as a term weighting scheme for information retrieval is a well-established technique in the field \cite{bird2009natural}. For keyword extraction, it complements POS-tag filtering naturally: TF-IDF handles corpus-level rarity while POS tags restrict terms to content-bearing word classes (nouns and adjectives), together yielding more human-readable keywords than either method alone.

Sentiment analysis of news text is notably harder than sentiment analysis of reviews or social media, because news writing is deliberately neutral and factual. \citet{hutto2014vader} introduced VADER (Valence Aware Dictionary and sEntiment Reasoner), a lexicon and rule-based method tuned on social media text that generalizes to other domains including news headlines. Using VADER to pseudo-label a training corpus is an established approach for building lightweight in-domain sentiment classifiers when gold-standard labels are unavailable.
